%% ma: L12-B08-0102
%% lop: 12
%% chuong: 2
%% bai: 8
%% dang: 12.8.3
%% loai: TN4
%% muc: NB
%% dap_an: "C"
%% nguon: "CĐ2.3.pdf giáo viên gửi 27/10/2026 (M07, Lê Minh Khôi – Nguyễn Anh Khôi, Chương II Chủ đề 3 Biểu thức tọa độ của các phép toán vectơ), Câu 12"
%% kien_thuc: "Bài 8 (độ dài vectơ, khoảng cách hai điểm)"
%% trang_thai: chinh-thuc
\begin{ex}
Trong không gian $Oxyz$, cho ba vectơ $\vec a=(-2;2;0)$, $\vec b=(2;2;0)$, $\vec c=(2;2;2)$. Độ dài của vectơ $\vec u=\vec a+\vec b+\vec c$ bằng
\choice{$6$}{$11$}{\True $2\sqrt{11}$}{$2\sqrt6$}
\loigiai{$\vec u=(2;6;2)$, $|\vec u|=\sqrt{4+36+4}=2\sqrt{11}$.}
\end{ex}
